
\begin{filecontents*}{twotimescales_philosophy_refs.bib}
@book{canguilhem1991normal,
  author    = {Canguilhem, Georges},
  title     = {The Normal and the Pathological},
  translator= {Fawcett, Carolyn R. and Cohen, Robert S.},
  publisher = {Zone Books},
  address   = {New York},
  year      = {1991}
}

@book{canguilhem2008knowledge,
  author    = {Canguilhem, Georges},
  title     = {Knowledge of Life},
  editor    = {Marrati, Paola and Meyers, Todd},
  translator= {Geroulanos, Stefanos and Ginsburg, Daniela},
  publisher = {Fordham University Press},
  address   = {New York},
  year      = {2008}
}

@incollection{sholl2020plastic,
  author    = {Sholl, Jonathan},
  title     = {Plastic, Variable, and Constructive: Renewing Canguilhem's Biological Normativity},
  booktitle = {Vital Norms: Canguilhem's The Normal and the Pathological in the Twenty-First Century},
  editor    = {M{\'e}thot, Pierre-Olivier},
  publisher = {Springer},
  address   = {Cham},
  year      = {2020},
  pages     = {255--291}
}

@article{engel1977newmodel,
  author  = {Engel, George L.},
  title   = {The Need for a New Medical Model: A Challenge for Biomedicine},
  journal = {Science},
  volume  = {196},
  number  = {4286},
  pages   = {129--136},
  year    = {1977},
  doi     = {10.1126/science.847460}
}

@article{thorell2024between,
  author  = {Thorell, Amanda},
  title   = {Health and Disease: Between Naturalism and Normativism},
  journal = {Philosophy of Science},
  volume  = {91},
  number  = {2},
  pages   = {449--467},
  year    = {2024},
  doi     = {10.1017/psa.2023.113}
}

@article{amoretti2022value,
  author  = {Amoretti, M. Cristina and Lalumera, Elisabetta},
  title   = {Wherein is the Concept of Disease Normative? From Weak Normativity to Value-Conscious Naturalism},
  journal = {Medicine, Health Care and Philosophy},
  volume  = {25},
  pages   = {47--60},
  year    = {2022},
  doi     = {10.1007/s11019-021-10048-x}
}

@book{varela1991embodied,
  author    = {Varela, Francisco J. and Thompson, Evan and Rosch, Eleanor},
  title     = {The Embodied Mind: Cognitive Science and Human Experience},
  publisher = {MIT Press},
  address   = {Cambridge, MA},
  year      = {1991}
}

@book{hollnagel2013resilient,
  editor    = {Hollnagel, Erik and Braithwaite, Jeffrey and Wears, Robert L.},
  title     = {Resilient Health Care},
  publisher = {Ashgate},
  address   = {Farnham},
  year      = {2013}
}

@article{rabinowitz2018planetary,
  author  = {Rabinowitz, Peter M. and Pappaioanou, Marguerite and Bardosh, Kevin L. and Conti, Lisa},
  title   = {A Planetary Vision for One Health},
  journal = {BMJ Global Health},
  volume  = {3},
  number  = {5},
  pages   = {e001137},
  year    = {2018},
  doi     = {10.1136/bmjgh-2018-001137}
}

@article{ruiz2023onehealth,
  author  = {Ruiz de Casta{\~n}eda, Rafael and Villers, Jennifer and Faerron Guzm{\'a}n, Carlos A. and Eslanloo, Turan and de Paula, Nicole and Machalaba, Catherine and others},
  title   = {One Health and Planetary Health Research: Leveraging Differences to Grow Together},
  journal = {The Lancet Planetary Health},
  volume  = {7},
  number  = {2},
  pages   = {e109--e111},
  year    = {2023},
  doi     = {10.1016/S2542-5196(23)00002-5}
}

@book{kooijman2010deb,
  author    = {Kooijman, Sebastiaan A. L. M.},
  title     = {Dynamic Energy Budget Theory for Metabolic Organisation},
  edition   = {3},
  publisher = {Cambridge University Press},
  address   = {Cambridge},
  year      = {2010},
  doi       = {10.1017/CBO9780511805400}
}

@article{marques2018amp,
  author  = {Marques, Gon{\c{c}}alo M. and Augustine, Starrlight and Lika, Konstadia and Pecquerie, Laure and Domingos, Tiago and Kooijman, Sebastiaan A. L. M.},
  title   = {The AmP Project: Comparing Species on the Basis of Dynamic Energy Budget Parameters},
  journal = {PLOS Computational Biology},
  volume  = {14},
  number  = {5},
  pages   = {e1006100},
  year    = {2018},
  doi     = {10.1371/journal.pcbi.1006100}
}

@article{olker2022ecotox,
  author  = {Olker, Jennifer H. and Elonen, Colleen M. and Pilli, Anne and Anderson, Arne and Kinziger, Brian and Erickson, Stephen and Skopinski, Michael and Pomplun, Anita and LaLone, Carlie A. and Russom, Christine L. and Hoff, Daniel J.},
  title   = {The ECOTOXicology Knowledgebase: A Curated Database of Ecologically Relevant Toxicity Tests to Support Environmental Research and Risk Assessment},
  journal = {Environmental Toxicology and Chemistry},
  volume  = {41},
  number  = {6},
  pages   = {1520--1539},
  year    = {2022},
  doi     = {10.1002/etc.5324}
}

@article{ward2018nitrate,
  author  = {Ward, Mary H. and Jones, Rena R. and Brender, Jean D. and de Kok, Theo M. and Weyer, Peter J. and Nolan, Bernard T. and Villanueva, Cristina M. and van Breda, Simone G.},
  title   = {Drinking Water Nitrate and Human Health: An Updated Review},
  journal = {International Journal of Environmental Research and Public Health},
  volume  = {15},
  number  = {7},
  pages   = {1557},
  year    = {2018},
  doi     = {10.3390/ijerph15071557}
}

@article{hosseini2021nitrate,
  author  = {Hosseini, Fatemeh and Majdi, Maryam and Naghshi, Sina and Sheikhhossein, Fatemeh and Djafarian, Kurosh and Shab-Bidar, Sakineh},
  title   = {Nitrate-nitrite Exposure through Drinking Water and Diet and Risk of Colorectal Cancer: A Systematic Review and Meta-analysis of Observational Studies},
  journal = {Clinical Nutrition},
  volume  = {40},
  number  = {5},
  pages   = {3073--3081},
  year    = {2021},
  doi     = {10.1016/j.clnu.2020.11.010}
}

@manual{who2022gdwq,
  author       = {{World Health Organization}},
  title        = {Guidelines for Drinking-water Quality: Fourth Edition Incorporating the First and Second Addenda},
  organization = {World Health Organization},
  address      = {Geneva},
  year         = {2022},
  url          = {https://www.who.int/publications/i/item/9789240045064}
}

@manual{epa2026npdwr,
  author       = {{U.S. Environmental Protection Agency}},
  title        = {National Primary Drinking Water Regulations},
  organization = {U.S. Environmental Protection Agency},
  year         = {2026},
  url          = {https://www.epa.gov/ground-water-and-drinking-water/national-primary-drinking-water-regulations}
}

@manual{epa2003cumulative,
  author       = {{U.S. Environmental Protection Agency}},
  title        = {Framework for Cumulative Risk Assessment},
  organization = {U.S. Environmental Protection Agency},
  number       = {EPA/600/P-02/001F},
  address      = {Washington, DC},
  year         = {2003},
  url          = {https://www.epa.gov/risk/framework-cumulative-risk-assessment}
}

@article{sexton2012cumulative,
  author  = {Sexton, Ken},
  title   = {Cumulative Risk Assessment: An Overview of Methodological Approaches for Evaluating Combined Health Effects from Exposure to Multiple Environmental Stressors},
  journal = {International Journal of Environmental Research and Public Health},
  volume  = {9},
  number  = {2},
  pages   = {370--390},
  year    = {2012},
  doi     = {10.3390/ijerph9020370}
}

@article{bopp2019mixtures,
  author  = {Bopp, Stephanie K. and Kienzler, Aude and Richarz, Andrea-Nicole and van der Linden, Sander C. and Paini, Alicia and Parissis, Nikolaos and Worth, Andrew P.},
  title   = {Regulatory Assessment and Risk Management of Chemical Mixtures: Challenges and Ways Forward},
  journal = {Critical Reviews in Toxicology},
  volume  = {49},
  number  = {2},
  pages   = {174--189},
  year    = {2019},
  doi     = {10.1080/10408444.2019.1579169}
}

@article{mathers2006burden,
  author    = {Mathers, Colin D. and Lopez, Alan D. and Murray, Christopher J. L.},
  title     = {The Burden of Disease and Mortality by Condition: Data, Methods, and Results for 2001},
  journal   = {Global Burden of Disease and Risk Factors},
  year      = {2006}
}
\end{filecontents*}

\documentclass[11pt]{article}

\usepackage[margin=1in]{geometry}
\usepackage{amsmath,amssymb,bm}
\usepackage[round,authoryear]{natbib}
\usepackage[hidelinks]{hyperref}
\usepackage{microtype}
\usepackage{enumitem}

\title{Two-Timescale Resilience and Adaptive Margin in Water--Health Assessment}
\author{Draft for discussion}
\date{\today}

\begin{document}
\maketitle

\section{Introduction}

Water-quality regulation has been most successful where harmful exposure can be linked to a well-characterised endpoint. Nitrate is a paradigmatic example. Drinking-water limits for nitrate and nitrite were historically justified by the prevention of infant methaemoglobinaemia: nitrate can be reduced to nitrite, nitrite can oxidise ferrous iron in haemoglobin, and the resulting methaemoglobin impairs oxygen transport. Contemporary regulatory values---for example the United States maximum contaminant level of 10 mg/L nitrate-nitrogen and World Health Organization drinking-water guideline values---retain this acute protective logic \citep{epa2026npdwr,who2022gdwq}. This logic should not be dismissed. It has made environmental hazards administratively visible, legally enforceable, and practically manageable.

The limitation is not that such thresholds are wrong, but that they represent only one evidentiary form: a contaminant-specific limit tied to a discrete adverse endpoint. Such reasoning is powerful when a causal chain is short, acute, and mechanistically well isolated. It becomes less adequate when water-mediated risk is chronic, cumulative, multi-stressor, spatially heterogeneous, persistent, bioaccumulative, or expressed as a gradual erosion of the target organism's capacity to respond before a diagnosable disease state appears. Cumulative risk assessment and mixture-risk research were developed precisely because single-stressor or single-threshold approaches cannot exhaust the structure of real exposure situations \citep{epa2003cumulative,sexton2012cumulative,bopp2019mixtures}. The challenge, therefore, is not simply to proclaim that water--health systems are complex. It is to develop tractable ways of representing what chronic complexity does to future vulnerability.

This paper begins from the claim that health is not adequately represented as mere absence of measured abnormality or conformity to a fixed normal state. Engel's critique of the reductionist biomedical model remains relevant because a purely biomedical model tends to define disease as deviation in measurable somatic variables while bracketing wider environmental, social, and behavioural relations \citep{engel1977newmodel}. Canguilhem's account of biological normativity gives this critique a specifically biological form: health is not mere conformity to a statistical norm, but the capacity of a living being to establish and re-establish viable norms in relation to a changing milieu \citep{canguilhem1991normal,canguilhem2008knowledge}. Contemporary readings of Canguilhem emphasise precisely this plastic, variable, organism--milieu character of biological normativity \citep{sholl2020plastic}.

In this paper, we use the term \emph{adaptive margin} for a formal proxy of one dimension of health: the remaining latitude to tolerate, absorb, recover from, or reorganise under perturbation. Adaptive margin is not identified with health as such. Rather, it represents a modelled component of vulnerability: the capacity that remains available when a target organism, population, life stage, or regulatory receptor encounters further stress. The central question therefore shifts from
\begin{quote}
  Has a contaminant-specific limit been exceeded?
\end{quote}
to
\begin{quote}
  How much less capable is this target system of absorbing what comes next?
\end{quote}

The conceptual structure can be stated compactly:
\begin{equation}
  \text{milieu pressure + exposure history}
  \rightarrow
  \text{narrowed normative latitude}
  \rightarrow
  \text{reduced restoring force}
  \rightarrow
  \text{amplified future burden}.
  \label{eq:intro_canguilhem_chain}
\end{equation}
Here, \emph{milieu pressure} refers to environmental and toxicological constraints; \emph{exposure history} refers to the persistence of past exposure in the present state of the system; \emph{normative latitude} refers to the capacity to sustain or re-establish viable functioning; and \emph{amplified burden} refers to the increased cost of a later perturbation under depleted background conditions.

To make this sequence operational, we translate it into a two-timescale model. On a slow timescale, background environmental pressures and retained internal burdens alter the state of a specified biological target by depleting adaptive margin. On a fast timescale, acute perturbations occur against this already-conditioned background. The same acute event can therefore produce different burdens depending on the prior state of the target system. Vulnerability is not represented only as exposure, hazard, or threshold exceedance, but as the amplification of acute burden caused by chronic depletion of response capacity.

Some minimal formalism is useful here because the conceptual move is also the modelling move. Let $A_t$ denote adaptive margin at background time $t$. In its simplest form,
\begin{equation}
  A_t = A_0 - \bm{\alpha}\cdot \mathbf{s}_t,
  \label{eq:intro_adaptive_margin}
\end{equation}
where $A_0$ is baseline adaptive margin, $\mathbf{s}_t$ is a vector of physiologically structured burdens, and $\bm{\alpha}$ weights how strongly those burdens deplete margin. This equation is not a neutral technical convenience. It commits the model to treating heterogeneous pressures as capable of being translated into a common loss of response latitude. It also makes clear that adaptive margin is target-specific: different organisms, populations, life stages, or receptors may begin with different margins and may lose them at different rates.

Adaptive margin then determines a restoring force, written $\lambda(A_t)$. The term \emph{restoring force} should not be read as a commitment to a fixed physiological normal state. In the present framework, it denotes the practical capacity to limit, absorb, or recover from displacement under a specified perturbation. If chronic background pressure reduces $A_t$, then the same acute event may impose a greater integrated burden. This motivates the amplification factor
\begin{equation}
  F_t = \frac{\lambda(A_0)}{\lambda(A_t)}.
  \label{eq:intro_amplification}
\end{equation}
When $F_t=1$, background conditions do not amplify the event relative to baseline. When $F_t>1$, chronic background pressure has made the same event more burdensome. This is the central object of the framework.

We call the framework \emph{TwoTimescaleResilience}. Its applied implementation follows a capacity--pressure--memory architecture. The capacity side is species- or target-specific physiological latitude: baseline adaptive margin, axis-specific depletion sensitivities, and bounds on restoring force. These quantities are informed by Dynamic Energy Budget (DEB) theory and by Add-my-Pet (AmP) records, which provide species-level DEB models, parameters, and underlying life-history data \citep{kooijman2010deb,marques2018amp}. The pressure side is stressor-specific evidence: no-effect and effect concentrations, endpoint categories, taxonomic groupings, and effect records. This layer is informed by the EPA ECOTOX Knowledgebase, a curated source of single-chemical toxicity test results across aquatic and terrestrial species \citep{olker2022ecotox}. The memory side separates ambient concentration from retained internal burden for persistent or bioaccumulative compounds.

In empirical use, the framework translates water-quality or environmental fields into effective exposure, effective exposure into toxicological pressure, toxicological pressure into burdens on DEB-informed physiological axes, physiological burdens into adaptive margin, adaptive margin into restoring force, and restoring force into acute-burden amplification. In compact form, the empirical chain is
\begin{equation}
  C_{j,t}\mapsto B_{j,t}\mapsto x_{j,t}\mapsto \mathbf{s}_{t}
  \mapsto A_t\mapsto\lambda(A_t)\mapsto F_t,
  \label{eq:intro_empirical_chain}
\end{equation}
where $C_{j,t}$ is ambient concentration of compound $j$, $B_{j,t}$ is internal burden, $x_{j,t}$ is active toxicological stress, $\mathbf{s}_t$ is DEB-axis burden, $A_t$ is adaptive margin, $\lambda(A_t)$ is restoring force, and $F_t$ is amplification. The full mathematical definition of each step is given below. The important point for the introduction is that each arrow is an interpretive commitment: it specifies how environmental presence becomes exposure, how exposure becomes pressure, how pressure depletes capacity, and how depleted capacity changes future vulnerability.

The nitrate case illustrates why this distinction matters without requiring us to claim that existing standards are simply wrong. The regulatory limit for nitrate was designed around infant methaemoglobinaemia, while subsequent epidemiological research has examined additional outcomes including colorectal cancer, thyroid disease, and neural tube defects \citep{ward2018nitrate}. Evidence for chronic outcomes is heterogeneous and, in some cases, contested; for example, a meta-analysis found dietary nitrate associated with colorectal cancer risk but did not find a statistically significant association for nitrate in drinking water alone \citep{hosseini2021nitrate}. This mixed evidence strengthens rather than weakens the modelling problem. Nitrate is not merely a toxicant with one endpoint; it is also part of nitrogen metabolism, diet, endogenous nitrosation, and physiological signalling. A threshold can protect against a known acute syndrome while still leaving unresolved how chronic exposure participates in system-level vulnerability.

TwoTimescaleResilience is therefore neither a rejection of regulatory thresholds nor a rejection of mechanistic toxicology. Thresholds remain necessary for enforceability and acute protection. Mechanistic toxicology remains necessary for identifying pathways and parameterising effects. But neither is sufficient for representing the upstream erosion of adaptive capacity under chronic, cumulative, persistent, bioaccumulative, and spatially heterogeneous stress. The contribution of this paper is to formalise that erosion as a two-timescale problem: chronic pressure narrows adaptive margin on a slow timescale, and this narrowed margin amplifies the burden of acute events on a fast timescale.

The remainder of the paper proceeds as follows. Section~\ref{sec:conceptual_scope} clarifies the conceptual scope and interpretive commitments of the framework. Section~\ref{sec:formal_translation} translates the capacity--pressure--memory structure into mathematical form. Subsequent sections describe the empirical parameterisation using AmP-derived capacity, ECOTOX-derived pressure, and compound memory; discuss uncertainty and aggregation; and consider how the framework can be used in water--health assessment without reducing health to a single threshold or vulnerability to a generic hazard index.

\section{Conceptual Scope and Interpretive Commitments}
\label{sec:conceptual_scope}

The framework is biologically parameterised, but its interpretation as vulnerability is value-conscious. Deciding that a loss of response capacity matters is not a purely descriptive act. Recent philosophy of medicine has moved beyond a simple naturalism-versus-normativism dichotomy and emphasises intermediate positions in which biological description and value-laden interpretation must both be made explicit \citep{thorell2024between,amoretti2022value}. The present framework follows that intermediate route. It uses biological quantities to estimate adaptive margin, but it does not pretend that the importance of lost margin can be inferred from biology alone.

This also clarifies the relation between adaptive margin and health. Adaptive margin is not health itself. It is a formal proxy for one dimension of health: the remaining latitude to sustain or re-establish viable relations under changing milieu conditions. In Canguilhem's terms, the organism is not merely defending a fixed normal state against external attack. Rather, it has a historically conditioned capacity to establish norms in relation to its milieu. Chronic environmental pressure and retained internal burden may narrow this capacity before a recognised disease state or endpoint failure appears.

The term \emph{restoring force} requires similar caution. In resilience engineering, resilience is often treated as the capacity of a system to sustain required operations under variable conditions \citep{hollnagel2013resilient}. This is useful, but in environmental health the point is not to ask organisms or communities to tolerate more pressure. The point is to reveal where adaptive margin is being consumed. In this paper, restoration does not mean return to a single fixed normal state. It means the modelled capacity to limit, absorb, recover from, or reorganise after perturbation.

The framework is also compatible with, but narrower than, One Health and Planetary Health. Those frameworks rightly emphasise the interdependence of human, animal, plant, ecosystem, and planetary health \citep{rabinowitz2018planetary,ruiz2023onehealth}. TwoTimescaleResilience operates at a more specific modelling level: target-specific physiological vulnerability under environmental pressure. It does not yet model trophic interactions, infectious transmission pathways, social institutions, environmental justice, ecosystem feedbacks, or planetary-scale governance. For human applications, adaptive margin should therefore not be read as an individual biological property alone; social and infrastructural conditions shape exposure, recovery capacity, and the distribution of future burden. For ecological applications, species-level amplification metrics should not be interpreted as rankings of ecological or moral value. They are conditional outputs of a particular parameterisation, endpoint selection, and exposure scenario.

The modelling unit is deliberately conventional: the named biological target used in empirical parameterisation, such as a species, population, life stage, or regulatory receptor. Thus, when the framework is parameterised for \emph{Homo sapiens}, \emph{Daphnia magna}, a fish species, or any other taxon, it does not explicitly model a host--symbiont assemblage or any alternative theory of biological individuality. Such relations may be implicit in empirical data used to estimate response, tolerance, or energetic parameters, but they are not the explicit object of the present model. TwoTimescaleResilience is a framework for background-conditioned physiological vulnerability, not a theory of what the biological individual ultimately is.

Finally, the framework should not be read as replacing existing modelling traditions. Cumulative risk assessment aggregates stressor-specific risks or hazard quotients; drinking-water cumulative toxicity scores combine co-occurring contaminants into sample-level indices; QMRA maps microbial dose into probabilities of infection and illness; PBPK/PBTK models translate external exposure into internal dose; TKTD, GUTS, and DEBtox models simulate time-dependent damage, survival, growth, or reproduction; AOP--DEB models connect mechanistic key events to energy-budget parameters; species sensitivity distributions infer protective concentrations from interspecies toxicity distributions; climate-vulnerability frameworks combine exposure, sensitivity, and adaptive capacity into an index; and early-warning approaches infer resilience loss from time-series statistics. TwoTimescaleResilience is adjacent to these approaches, but its narrower contribution is to formalise how chronic background burden changes the restoring force available during later acute perturbations.

\section{Translating Capacity--Pressure--Memory into Mathematical Form}
\label{sec:formal_translation}

The mathematical formulation does not simply operationalise the conceptual sequence; it interprets it. Each modelling step---defining exposure, assigning stressor pressure, aggregating physiological burden, representing adaptive margin, and computing amplification---commits the framework to a particular account of how organisms relate to their milieu. This section therefore presents the equations as modelling commitments rather than as neutral transformations.

Let $g$ denote a spatial unit, $t$ a slow background time index, and $q$ a species, population, life stage, or target biological profile. Let
\begin{equation}
  \mathbf{R}_{g,t} = (R_{1,g,t},\ldots,R_{J,g,t})^\top
\end{equation}
be a vector of raw environmental stressors. Raw environmental presence is not yet exposure. A target-specific exposure filter maps raw environmental variables into effective exposures,
\begin{equation}
  \mathbf{e}^{(q)}_{g,t} = \mathbf{H}^{(q)}\mathbf{R}_{g,t},
  \label{eq:exposure_filter}
\end{equation}
where $\mathbf{H}^{(q)}$ encodes habitat use, contact, uptake, physiology, behaviour, or other exposure modifiers. This step expresses the idea that the milieu is target-relative: the relevant environment is the environment as encountered by a specified living system.

Effective exposure is then mapped into biological modes of action,
\begin{equation}
  \mathbf{m}^{(q)}_{g,t} = \mathbf{U}^{(q)}\mathbf{e}^{(q)}_{g,t},
  \label{eq:moa_mapping}
\end{equation}
and from modes of action into DEB-informed physiological axes,
\begin{equation}
  \mathbf{s}^{(q)}_{g,t} = \mathbf{W}^{(q)}\mathbf{m}^{(q)}_{g,t}.
  \label{eq:deb_axis_mapping}
\end{equation}
The entries of $\mathbf{s}$ are burdens on assimilation, maintenance, growth, and reproduction. This step is both biological and interpretive: it specifies which dimensions of viable functioning the model recognises as relevant to adaptive margin.

For empirical chemical stressors, ECOTOX-derived effect summaries define active stress. For contaminant $j$, a simple scaled burden is
\begin{equation}
  x_{j,t} = \max\left(0,\frac{E_{j,t}-NOEC_j}{EC50_j-NOEC_j}\right),
  \label{eq:ecotox_scaling}
\end{equation}
where $E_{j,t}$ is the effective concentration used for ECOTOX comparison, $NOEC_j$ is a taxon- and endpoint-relevant no-observed-effect concentration, and $EC50_j$ is a corresponding median effect concentration. In stateless calculations, $E_{j,t}=C_{j,t}$, the ambient concentration. In stateful calculations for persistent or bioaccumulative compounds, $E_{j,t}=B_{j,t}$, an internal burden state. ECOTOX records are therefore treated as curated experimental evidence of stressor pressure, not as context-free measures of harm. Routing effect codes to DEB axes is a modelling decision that must be reported transparently.

Compound memory is represented by
\begin{equation}
  B_{j,t}=\rho_j B_{j,t-1}+(1-\rho_j)K_j C_{j,t},
  \label{eq:compound_memory}
\end{equation}
where $C_{j,t}$ is ambient concentration, $\rho_j\in[0,1)$ is monthly retention or carryover, and $K_j\geq 0$ is a bioaccumulation or internal magnification factor. The limiting case $\rho_j=0$ and $K_j=1$ recovers the stateless relation $B_{j,t}=C_{j,t}$. If $\rho_j>0$ and $K_j=1$, the model represents carryover without internal magnification. If $\rho_j>0$ and $K_j>1$, the model represents both carryover and magnified internal burden. Chemical memory $B_{j,t}$ is distinct from any future organism-level physiological condition memory $Z_t$.

Adaptive margin is then defined as
\begin{equation}
  A^{(q)}_{g,t}
  = A^{(q)}_0 + \omega_Z Z^{(q)}_{g,t}
  - \bm{\alpha}^{(q)}\cdot \mathbf{s}^{(q)}_{g,t},
  \label{eq:A}
\end{equation}
where $A^{(q)}_0$ is baseline margin, $Z$ is an optional condition or reserve buffer, $\omega_Z$ scales the contribution of that buffer, and $\bm{\alpha}^{(q)}$ weights depletion associated with physiological-axis burden. In the current empirical implementation, $Z_t$ is not the chemical memory state; chemical carryover is handled upstream through $B_{j,t}$. The quantities $A_0^{(q)}$ and $\bm{\alpha}^{(q)}$ are capacity-side parameters: they express how much physiological margin the target has and how strongly burdens on assimilation, maintenance, growth, and reproduction erode that margin.

Adaptive margin determines a restoring force $\lambda$. A simple bounded form is
\begin{equation}
  \lambda(A) = \lambda_{\min} + (\lambda_{\max}-\lambda_{\min})
  \frac{A_+}{K_A + A_+},
  \qquad A_+ = \max(A,0),
  \label{eq:lambda}
\end{equation}
where $\lambda_{\min}$ and $\lambda_{\max}$ define lower and upper recovery bounds and $K_A$ is a half-saturation constant. The dependence of $\lambda$ on $A$ links recovery capacity to background-conditioned vulnerability: the system may lose the margin needed to absorb further disturbance before crossing into a narrower mode of functioning.

Let $C_{\mathrm{event}}(\tau)$ be an acute event profile on the fast time variable $\tau$, and let $y(\tau)$ denote the displacement or burden induced by that event. For a fixed background state $A_{g,t}^{(q)}$, consider
\begin{equation}
  \frac{dy}{d\tau} = -\lambda(A^{(q)}_{g,t})y
  + \beta C_{\mathrm{event}}(\tau),
  \qquad y(0)=0.
  \label{eq:pulse}
\end{equation}
Integrating Equation~\eqref{eq:pulse} over time gives
\begin{equation}
  \int_0^\infty y(\tau)\,d\tau
  =
  \frac{\beta}{\lambda(A^{(q)}_{g,t})}
  \int_0^\infty C_{\mathrm{event}}(\tau)\,d\tau,
  \label{eq:auc}
\end{equation}
provided $C_{\mathrm{event}}$ is integrable and $\lambda>0$. Thus, for identical acute event cost, the integrated displacement burden is inversely proportional to the background-conditioned restoring force. This yields the amplification factor
\begin{equation}
  F^{(q)}_{g,t}
  =
  \frac{\lambda(A^{(q)}_0)}
       {\lambda(A^{(q)}_{g,t})}.
  \label{eq:F}
\end{equation}
When $F=1$, background conditions do not amplify the event relative to baseline. When $F>1$, chronic background stress has made the same event more burdensome.

Taken together, the formal chain is
\begin{equation}
  C_{j,t}\mapsto B_{j,t}\mapsto x_{j,t}
  \mapsto \mathbf{s}_{t}
  \mapsto A_t
  \mapsto \lambda(A_t)
  \mapsto F_t.
  \label{eq:memory_chain}
\end{equation}
This chain is the operational form of the capacity--pressure--memory structure. AmP-derived DEB parameters describe target capacity; ECOTOX-derived effect summaries describe environmental pressure; compound memory separates ambient exposure from retained internal burden; and the amplification factor expresses the consequence of narrowed response capacity for future acute perturbations.

\section{Empirical Parameterisation}
\label{sec:empirical_parameterisation}

The applied implementation assigns the three conceptual components of the framework to distinct empirical layers. The capacity layer is informed by DEB theory and Add-my-Pet records. These provide species-level life-history and energetic parameters that can be translated into baseline adaptive margin, axis-specific depletion sensitivities, and lower and upper bounds on restoring force \citep{kooijman2010deb,marques2018amp}. The pressure layer is informed by ECOTOX-derived toxicity records, including taxonomic groupings, endpoint categories, no-effect concentrations, and effect concentrations \citep{olker2022ecotox}. The memory layer is compound-specific and separates ambient concentration from retained internal burden using retention and internal magnification parameters.

This division is useful because it prevents three distinct questions from being collapsed into one. First, what capacity does the target have? Second, what pressure is exerted by the stressor under the relevant exposure scenario? Third, how much of past exposure remains active in the present? A generic hazard index can combine environmental measurements, but the present framework asks how those measurements alter the response capacity of a specified target system.

Current runtime aggregation is axis-wise after endpoint routing. ECOTOX effect codes are mapped to DEB-informed axes, active stress is computed relative to no-effect and effect concentrations, and the resulting burdens are aggregated by axis before being passed to the adaptive-margin equation. This is intentionally transparent but incomplete. Formal concentration addition, independent action, and other mixture-toxicity logics are reserved for later extensions. The present paper therefore treats mixture aggregation as a reported modelling assumption rather than as a solved toxicological problem.

\section{Implications for Water--Health Assessment}
\label{sec:implications}

The central implication is that water--health assessment should distinguish exceedance from vulnerability. Exceedance asks whether a contaminant concentration passes a regulatory or toxicological limit. Vulnerability, in the sense used here, asks whether chronic pressure has reduced the target's capacity to withstand later perturbation. These are related but not equivalent. A target can remain below a regulatory threshold while still losing adaptive margin under cumulative or persistent pressure. Conversely, a threshold exceedance may be important for acute protection without fully describing the longer-term state of response capacity.

This distinction is particularly relevant for persistent and bioaccumulative compounds, but it is not limited to them. Compound memory makes the issue mathematically explicit: internal burden may remain elevated after ambient concentration declines. More generally, environmental histories may condition the present in ways that are invisible if assessment is based only on contemporaneous concentration or single-sample exceedance. TwoTimescaleResilience therefore provides a way to express background-conditioned vulnerability without abandoning the practical importance of regulatory thresholds.

The framework also changes the interpretation of acute events. An acute pulse is not evaluated only by its own magnitude. Its burden depends on the state into which it arrives. The same event may be relatively absorbable under high adaptive margin and more consequential under depleted margin. In this sense, the amplification factor $F$ is not a new disease endpoint. It is a measure of how much chronic background conditions have altered the cost of what comes next.

\section{Limitations and Future Work}
\label{sec:limitations}

Several limitations follow directly from the modelling choices above. First, adaptive margin is scalar in the present formulation, even though biological normativity is multidimensional. The scalar representation is useful for operational vulnerability assessment, but it inevitably compresses heterogeneous forms of capacity into a common metric. Second, the current implementation treats DEB-axis burdens as additive after endpoint routing. This should be replaced or supplemented by more explicit mixture-toxicity models where data permit. Third, chemical memory $B_t$ is included, but organism-level physiological condition memory $Z_t$ remains optional and is not yet the central empirical state variable. Fourth, social, infrastructural, behavioural, and institutional determinants of exposure and recovery are not explicitly modelled, even though they are crucial for human water--health applications.

Future work should therefore proceed in stages. The first priority is diagnostic validation of monthly memory behaviour across species, compounds, and exposure scenarios. The second is more explicit mixture handling, distinguishing cases where concentration addition, independent action, or bounded axis aggregation is appropriate. The third is the careful introduction of physiological condition memory, kept conceptually distinct from chemical internal burden. Only after these steps should the framework be connected to large-scale monthly water-quality rasters or used for broad spatial comparison.

\section{Conclusion}
\label{sec:conclusion}

This paper has proposed TwoTimescaleResilience as a framework for representing background-conditioned vulnerability in water--health assessment. Its starting point is the Canguilhemian idea that health is not mere conformity to a fixed normal state, but a capacity to establish and re-establish viable norms in relation to a changing milieu. The modelling translation of this idea is capacity--pressure--memory: environmental pressure and retained burden can narrow adaptive margin, reduced adaptive margin weakens restoring force, and weakened restoring force amplifies the burden of future acute perturbations.

The framework does not reject thresholds, mechanistic toxicology, or existing risk-assessment traditions. Rather, it addresses a narrower gap: how to represent the erosion of response capacity before or alongside discrete endpoint failure. By distinguishing ambient concentration from internal burden, toxicological pressure from physiological capacity, and acute event magnitude from background-conditioned burden, TwoTimescaleResilience offers a tractable way to ask not only whether a limit has been exceeded, but how much capacity remains to absorb what comes next.

\bibliographystyle{plainnat}
\bibliography{twotimescales_philosophy_refs}

\end{document}
